\documentclass[12pt]{ximera}
\title{Exam 3}
\begin{document}
\begin{abstract}
Exam 3, covering Chapter 3: matching fair division methods to continuous and discrete
situations, the method of markers on a candy array, sealed bids applied to household
chores (where every value is negative), and a lone-divider sundae that ends in a
standoff.
\end{abstract}
\maketitle

\section*{Exam 3}

A useful formula throughout:
\[
\text{value of a mixed share} = (\text{proportion of asset 1})\cdot(\text{value of asset 1})
+ (\text{proportion of asset 2})\cdot(\text{value of asset 2}).
\]

\begin{problem}
Identify which fair division methods could possibly apply to each situation.

\begin{prompt}
\textbf{(a)} Three siblings need to fairly divide the land from their grandpa's ranch
among themselves.

\begin{selectAll}
\choice[correct]{Lone-Divider Method}
\choice{Sealed Bids Method}
\choice[correct]{Lone-Chooser Method}
\choice{Method of Markers}
\end{selectAll}

\begin{hint}
Is land continuous or discrete? Which two of these methods are built for each kind of
asset?
\end{hint}

\begin{feedback}
Land is a \emph{continuous} asset --- it can be divided into pieces of any size --- so
the continuous methods apply: \textbf{Lone-Divider} and \textbf{Lone-Chooser}.

Sealed Bids and the Method of Markers are both designed for \emph{discrete} assets:
collections of indivisible objects.
\end{feedback}
\end{prompt}

\begin{prompt}
\textbf{(b)} Four business partners are separating from a shared business and need to
divide $550$ owned apartments and other properties among them.

\begin{selectAll}
\choice{Lone-Divider Method}
\choice[correct]{Sealed Bids Method}
\choice{Lone-Chooser Method}
\choice[correct]{Method of Markers}
\end{selectAll}

\begin{feedback}
Individual apartments and properties are \emph{discrete} --- you cannot hand someone
$40\%$ of an apartment building in this setting --- so the discrete methods apply:
\textbf{Sealed Bids} and the \textbf{Method of Markers}.

With $550$ items the Method of Markers is especially well suited, since it needs many
items of comparatively small individual value in order to work well. Sealed Bids works
too, though it requires the partners to have enough cash on hand to settle up.
\end{feedback}
\end{prompt}
\end{problem}

\begin{problem}
In which of the following fair division methods is it possible for \emph{every} player
to come out with \emph{greater} than their minimum fair share?

\begin{selectAll}
\choice{Lone-Divider Method}
\choice[correct]{Sealed Bids Method}
\choice{Lone-Chooser Method}
\choice[correct]{Method of Markers}
\end{selectAll}

\begin{hint}
In the divider-based methods, think about what the \emph{divider} ends up with. Can the
divider ever do better than exactly a fair share?
\end{hint}

\begin{feedback}
\textbf{Sealed Bids} and the \textbf{Method of Markers}.

In \emph{Sealed Bids}, each player is first brought to exactly their fair share, and
then the surplus --- which exists whenever the players value the items differently ---
is split equally among them. Every player ends strictly better off than a fair share.
Homework 8 is a worked example: all four roommates finished \$52.50 ahead.

In the \emph{Method of Markers}, every player receives a segment they value at a fair
share or more, and the leftover items are then distributed among them, pushing
everyone above the line.

The two divider-based methods cannot do this, because of the \textbf{divider}. The
divider cuts the asset into pieces they consider exactly equal, so whichever piece they
end up with is worth exactly $\frac1N$ to them --- never more. The choosers may do
better, but the divider cannot, so ``every player'' fails.
\end{feedback}
\end{problem}

\begin{problem}
In the following discrete system with four players, we use the Method of Markers. The
$21$ candies are laid out in an array, and each player's three markers are shown
beneath it.

\begin{center}
\includegraphics[width=\textwidth]{candies.jpg}
\end{center}

The markers fall as follows:
\[
\begin{array}{ll}
\text{Player A}: & 4,\ 9,\ 17\\
\text{Player B}: & 5,\ 11,\ 16\\
\text{Player C}: & 6,\ 10,\ 15\\
\text{Player D}: & 6,\ 12,\ 16
\end{array}
\]

\begin{prompt}
Which items will go to each player, and which items will be surplus?

Round 1 goes to player $\answer{A}$, who takes items 1--$\answer{4}$.

\smallskip
Round 2 goes to player $\answer{C}$, who takes items 5--$\answer{10}$.

\smallskip
Round 3 is a tie between players B and D at marker $\answer{16}$. The winner of the
tiebreak takes items 11--16, and the remaining player takes items
$\answer{17}$--21.

\smallskip
The surplus consists of $\answer{0}$ items.

\begin{hint}
In each round, compare only the relevant markers (first markers in round 1, second in
round 2, and so on) among the players not yet served, and find the leftmost one lying
past the items already distributed.
\end{hint}

\begin{feedback}
\textbf{Round 1.} First markers are $A_1=4$, $B_1=5$, $C_1=6$, $D_1=6$. A's is
leftmost, so \textbf{A takes items 1--4} and leaves.

\textbf{Round 2.} Among B, C, and D, the second markers are $B_2=11$, $C_2=10$,
$D_2=12$. C's is leftmost, so \textbf{C takes items 5--10}.

\textbf{Round 3.} Only B and D remain, and their third markers are both at $16$ --- a
genuine tie, which is resolved by a coin flip or other random tiebreak. The winner
takes items \textbf{11--16}, and the other player, being last, takes everything past
their final marker: items \textbf{17--21}.

\textbf{Surplus: none.} The four segments $1$--$4$, $5$--$10$, $11$--$16$, and
$17$--$21$ cover all $21$ candies with nothing left over.

That is unusual --- the Method of Markers normally leaves a surplus --- and it happens
here precisely because $B_3$ and $D_3$ landed on the same marker. Had those two markers
been even one position apart, the candies between them would have been surplus.
\end{feedback}
\end{prompt}
\end{problem}

\begin{problem}
Jake, Adam, and Joel want to share a $12$-inch ultra-sundae made up of half vanilla and
half chocolate ice cream, using the lone-divider method. Adam will be the divider.
\begin{itemize}
\item Adam values the chocolate three times as much as the vanilla;
\item Jake values the vanilla twice as much as the chocolate;
\item Joel has no preference of one over the other.
\end{itemize}

\begin{prompt}
\textbf{(a)} Complete the chart with each player's value system for this sundae, as
fractions of the whole.

Adam --- Vanilla: $\answer{1/4}$, \quad Chocolate: $\answer{3/4}$

\smallskip
Jake --- Vanilla: $\answer{2/3}$, \quad Chocolate: $\answer{1/3}$

\smallskip
Joel --- Vanilla: $\answer{1/2}$, \quad Chocolate: $\answer{1/2}$

\begin{hint}
``Values chocolate three times as much as vanilla'' means the two halves are in a
$1:3$ ratio of value. The two fractions must add to $1$.
\end{hint}

\begin{feedback}
\textbf{Adam}: chocolate is worth $3$ times vanilla, a $1:3$ split, so vanilla is
$\frac14$ and chocolate $\frac34$.

\textbf{Jake}: vanilla is worth twice the chocolate, a $2:1$ split, so vanilla is
$\frac23$ and chocolate $\frac13$.

\textbf{Joel}: no preference, so $\frac12$ each.
\end{feedback}
\end{prompt}

\begin{prompt}
\textbf{(b)} Describe the three shares Adam will make, as a portion of Vanilla (V) plus
a portion of Chocolate (C).

$s_1$ = all of V $+$ $\answer{1/9}$ of C \qquad
$s_2 = \answer{4/9}$ of C \qquad
$s_3 = \answer{4/9}$ of C

\begin{hint}
Adam must make three shares each worth $\frac13$ to him. All of the vanilla is worth
only $\frac14$ --- less than $\frac13$ --- so the first cut has to reach into the
chocolate.
\end{hint}

\begin{feedback}
Adam needs three shares worth $\frac13$ each in his own values, where V $= \frac14$ and
C $= \frac34$.

All of the vanilla is worth only $\frac14 < \frac13$, so $s_1$ must include all of V
plus some chocolate. The shortfall is
\[ \tfrac13 - \tfrac14 = \tfrac{1}{12}, \]
and since all of C is worth $\frac34$, the fraction of C needed is
\[ \frac{1/12}{3/4} = \frac19. \]
So $s_1 =$ all of V $+\ \frac19$ of C.

That leaves $\frac89$ of the chocolate, worth $\frac89 \cdot \frac34 = \frac23$ to
Adam, which splits into two shares of $\frac13$ each:
\[ s_2 = s_3 = \tfrac49 \text{ of C}. \]
Check: $\frac19+\frac49+\frac49 = 1$, so all the chocolate is accounted for. \checkmark

(On the $12$-inch sundae, that is a cut $\frac23$ of an inch into the chocolate, then
two equal $\frac{8}{3}$-inch pieces.)
\end{feedback}
\end{prompt}

\begin{prompt}
\textbf{(c)} What fraction of the whole does each player consider each share to be
worth? Give your answers as fractions.

Adam --- $s_1: \answer{1/3}$, \quad $s_2: \answer{1/3}$, \quad $s_3: \answer{1/3}$

\smallskip
Jake --- $s_1: \answer{19/27}$, \quad $s_2: \answer{4/27}$, \quad $s_3: \answer{4/27}$

\smallskip
Joel --- $s_1: \answer{5/9}$, \quad $s_2: \answer{2/9}$, \quad $s_3: \answer{2/9}$

\begin{feedback}
Apply the mixed-share formula with each player's own values.

\textbf{Adam} ($V=\frac14$, $C=\frac34$):
\[ s_1 = \tfrac14 + \tfrac19\cdot\tfrac34 = \tfrac14+\tfrac{1}{12} = \tfrac13,
\qquad s_2 = s_3 = \tfrac49\cdot\tfrac34 = \tfrac13. \]
As the divider, all three are equal for him --- as they must be.

\textbf{Jake} ($V=\frac23$, $C=\frac13$):
\[ s_1 = \tfrac23 + \tfrac19\cdot\tfrac13 = \tfrac{18}{27}+\tfrac{1}{27} = \tfrac{19}{27} \approx 70.4\%, \]
\[ s_2 = s_3 = \tfrac49\cdot\tfrac13 = \tfrac{4}{27} \approx 14.8\%. \]

\textbf{Joel} ($V=\frac12$, $C=\frac12$):
\[ s_1 = \tfrac12 + \tfrac19\cdot\tfrac12 = \tfrac{10}{18} = \tfrac59 \approx 55.6\%,
\qquad s_2 = s_3 = \tfrac49\cdot\tfrac12 = \tfrac29 \approx 22.2\%. \]

Each row totals $1$. \checkmark

A fair share is $\frac13 \approx 33.3\%$, so the bid lists are: Adam $\set{s_1,s_2,s_3}$,
Jake $\set{s_1}$, Joel $\set{s_1}$.
\end{feedback}
\end{prompt}

\begin{prompt}
\textbf{(d)} What is one fair division of the sundae among the three friends?

\begin{freeResponse}
\end{freeResponse}

\begin{hint}
Look at the bid lists from part (c). Can Jake and Joel both be satisfied with whole
shares?
\end{hint}

\begin{feedback}
\textbf{There is no fair division using the shares as cut} --- this is a standoff. Both
Jake and Joel consider \emph{only} $s_1$ to be worth a fair share, and there is just one
$s_1$.

The standard resolution is to settle the uncontested player first and re-divide the
rest:
\begin{enumerate}
\item Give $s_2$ (or $s_3$) to \textbf{Adam}. As the divider he values every share at
      exactly $\frac13$, so this is fair for him.
\item Recombine the remaining two shares --- $s_1$ and $s_3$ --- into a single lot, and
      have \textbf{Jake and Joel} divide it between themselves using divider-chooser.
\end{enumerate}

Check that this works. Jake values $s_1 + s_3$ at
$\frac{19}{27}+\frac{4}{27} = \frac{23}{27} \approx 85.2\%$ of the sundae, and Joel
values it at $\frac59 + \frac29 = \frac79 \approx 77.8\%$. Divider-chooser guarantees
each of them at least half of what they think the lot is worth --- at least $42.6\%$ for
Jake and at least $38.9\%$ for Joel. Both clear the $33.3\%$ they need, so every player
ends up with a fair share.
\end{feedback}
\end{prompt}
\end{problem}

\begin{problem}
Jake and their partner Jessie live together and are using the method of sealed bids to
split responsibilities for chores. They assign \emph{negative} values to each chore,
based on how much they would be willing to pay the other to do it instead. A
less-negative value means the task is less repulsive to that person.

\begin{center}
\begin{tabular}{|r|c|c|}\hline
 & \textbf{Jake} & \textbf{Jessie} \\\hline
Run the dishwasher & $-\$2$ & $-\$3$ \\
Do the laundry & $-\$3$ & $-\$5$ \\
Plan meals for the week & $-\$4$ & $-\$9$ \\
Shop for groceries & $-\$2$ & $-\$8$ \\
Take out the trash & $-\$3$ & $-\$2$ \\
Clean the kitchen & $-\$5$ & $-\$6$ \\
Clean the bathroom & $-\$8$ & $-\$4$ \\
Water the houseplants & $-\$7$ & $-\$1$ \\\hline
\end{tabular}
\end{center}

\begin{prompt}
\textbf{(a)} Based on the listed bids, who will be assigned each task?

Dishwasher: $\answer{Jake}$ \quad Laundry: $\answer{Jake}$ \quad
Meal planning: $\answer{Jake}$ \quad Groceries: $\answer{Jake}$

\smallskip
Trash: $\answer{Jessie}$ \quad Kitchen: $\answer{Jake}$ \quad
Bathroom: $\answer{Jessie}$ \quad Houseplants: $\answer{Jessie}$

\begin{hint}
Each item still goes to the \emph{highest} bidder. With negative numbers, the highest
bid is the least negative one --- the person who minds the chore least.
\end{hint}

\begin{feedback}
Each chore goes to whoever values it highest, which with negative numbers means
whoever minds it least:
\[
\begin{array}{lll}
\text{Dishwasher} & -2 > -3 & \text{Jake}\\
\text{Laundry} & -3 > -5 & \text{Jake}\\
\text{Meal planning} & -4 > -9 & \text{Jake}\\
\text{Groceries} & -2 > -8 & \text{Jake}\\
\text{Trash} & -2 > -3 & \text{Jessie}\\
\text{Kitchen} & -5 > -6 & \text{Jake}\\
\text{Bathroom} & -4 > -8 & \text{Jessie}\\
\text{Houseplants} & -1 > -7 & \text{Jessie}
\end{array}
\]
\end{feedback}
\end{prompt}

\begin{prompt}
\textbf{(b)} What are Jake's and Jessie's fair shares in this system? (These will be
negative.)

Jake: $\answer{-17}$ \qquad Jessie: $\answer{-19}$

\begin{feedback}
Jake's total valuation of all the chores is
\[ -2-3-4-2-3-5-8-7 = -\$34, \]
so his fair share is $-34/2 = -\$17$.

Jessie's total is
\[ -3-5-9-8-2-6-4-1 = -\$38, \]
so her fair share is $-38/2 = -\$19$.

A fair share here means each should end up doing chores they find no worse than that
amount of unpleasantness.
\end{feedback}
\end{prompt}

\begin{prompt}
\textbf{(c)} Compare part (b) with each partner's perceived value of the tasks they've
been assigned. What do they each owe, or what are they owed, in the initial settlement?

Jake's allotment value: $\answer{-16}$ \qquad Jessie's allotment value: $\answer{-7}$

\smallskip
Initial settlement --- Jake pays $\$\answer{1}$ and Jessie pays $\$\answer{12}$.

\begin{feedback}
Jake takes the dishwasher, laundry, meal planning, groceries, and kitchen:
\[ -2-3-4-2-5 = -\$16. \]
Jessie takes the trash, bathroom, and houseplants:
\[ -2-4-1 = -\$7. \]

Compare each with the fair share:
\begin{itemize}
\item Jake: $-16$ versus a fair share of $-17$. His chore load is $\$1$ \emph{less}
      unpleasant than fair, so he pays $\$1$ into the system.
\item Jessie: $-7$ versus a fair share of $-19$. Her load is $\$12$ less unpleasant
      than fair, so she pays $\$12$ in.
\end{itemize}
Both partners come out ahead of their fair share on chores alone, so both pay in ---
which is what creates the surplus.
\end{feedback}
\end{prompt}

\begin{prompt}
\textbf{(d)} What is the surplus after the initial settlement?

$\$\answer{13}$ \qquad and each partner's half is $\$\answer{6.50}$

\begin{feedback}
Jake pays \$1 and Jessie pays \$12, so the system holds
\[ \$1 + \$12 = \$13, \]
and nothing is paid out. Split evenly, that is $\$13/2 = \$6.50$ each.
\end{feedback}
\end{prompt}

\begin{prompt}
\textbf{(e)} After splitting the surplus equally, how much money is each partner making
or paying? Enter a positive number for money received, negative for money paid.

Jake: $\answer{5.50}$ \qquad Jessie: $\answer{-5.50}$

\begin{feedback}
Subtract each partner's \$6.50 share of the surplus from what they owed:
\begin{itemize}
\item Jake: paid \$1, receives \$6.50 back, so he nets $+\$5.50$ --- he \emph{receives}
      \$5.50.
\item Jessie: paid \$12, receives \$6.50 back, so she nets $-\$5.50$ --- she
      \emph{pays} \$5.50.
\end{itemize}

So Jessie pays Jake \$5.50 a week. The final arrangement:
\begin{itemize}
\item \textbf{Jake} runs the dishwasher, does laundry, plans meals, shops for
      groceries, and cleans the kitchen, and receives \$5.50.
\item \textbf{Jessie} takes out the trash, cleans the bathroom, waters the houseplants,
      and pays \$5.50.
\end{itemize}
Check: Jake's position is $-16 + 5.50 = -\$10.50$ against a fair share of $-\$17$, so he
is \$6.50 ahead. Jessie's is $-7 - 5.50 = -\$12.50$ against $-\$19$, also \$6.50 ahead.
\checkmark
\end{feedback}
\end{prompt}
\end{problem}

\end{document}
